\textbf{Phần 1. Trắc nghiệm lựa chọn}
\setcounterpageref{ex}{0}

\begin{ex} %[4D1N02]
Kết quả của phép cộng $45\,236 + 23\,142$ là
\choice
{$68\,377$}
{\True $68\,378$}
{$68\,478$}
{$67\,378$}
\loigiai{
$45\,236 + 23\,142 = 68\,378$.\\
\textbf{Đáp án của câu hỏi là B}
}
\end{ex}

\begin{ex} %[4D1N02]
Kết quả của phép trừ $78\,560 - 34\,215$ là
\choice
{\True $44\,345$}
{$44\,355$}
{$44\,335$}
{$43\,345$}
\loigiai{
$78\,560 - 34\,215 = 44\,345$.\\
\textbf{Đáp án của câu hỏi là A}
}
\end{ex}

\begin{ex} %[4D1H02]
Kết quả của phép nhân $12\,304 \times 5$ là
\choice
{$61\,500$}
{$60\,520$}
{$61\,525$}
{\True $61\,520$}
\loigiai{
$12\,304 \times 5 = 61\,520$.\\
\textbf{Đáp án của câu hỏi là D}
}
\end{ex}

\begin{ex} %[4D1H02]
Kết quả của phép chia $84\,672 : 4$ là
\choice
{$21\,167$}
{$21\,158$}
{\True $21\,168$}
{$20\,168$}
\loigiai{
$84\,672 : 4 = 21\,168$.\\
\textbf{Đáp án của câu hỏi là C}
}
\end{ex}

\begin{ex} %[4D1H02]
Giá trị của biểu thức $15\,000 + 25\,000 : 5$ là
\choice
{$8000$}
{$25\,000$}
{$18\,000$}
{\True $20\,000$}
\loigiai{
$15\,000 + 5000 = 20\,000$.\\
\textbf{Đáp án của câu hỏi là D}
}
\end{ex}

\begin{ex} %[4D1H02]
Tìm $x$ biết: $x - 12\,500 = 34\,200$.
\choice
{\True $46\,700$}
{$21\,700$}
{$46\,500$}
{$47\,700$}
\loigiai{
$x = 34\,200 + 12\,500 = 46\,700$.\\
\textbf{Đáp án của câu hỏi là A}
}
\end{ex}

\begin{ex} %[4D1V02]
Một cửa hàng lương thực nhập về $2$ chuyến xe chở gạo, chuyến thứ nhất chở $14\,500\text{ kg}$ gạo, chuyến thứ hai chở gấp đôi chuyến thứ nhất. Cả hai chuyến chở được tất cả số ki-lô-gam gạo là
\choice
{$29\,000\text{ kg}$}
{$42\,500\text{ kg}$}
{\True $43\,500\text{ kg}$}
{$45\,000\text{ kg}$}
\loigiai{
Chuyến thứ hai chở: $14\,500 \times 2 = 29\,000\text{ kg}$.\\
Cả hai chuyến chở: $14\,500 + 29\,000 = 43\,500\text{ kg}$.\\
\textbf{Đáp án của câu hỏi là C}
}
\end{ex}

\begin{ex} %[4D1V02]
Tính giá trị biểu thức: $(48\,600 - 18\,600) \times 3 : 6$.
\choice
{$12\,000$}
{\True $15\,000$}
{$18\,000$}
{$20\,000$}
\loigiai{
$30\,000 \times 3 : 6 = 90\,000 : 6 = 15\,000$.\\
\textbf{Đáp án của câu hỏi là B}
}
\end{ex}

\begin{ex} %[4D1V02]
Một công ty may trong 5 ngày may được $12\,500$ chiếc áo sơ mi. Hỏi trong 8 ngày với cùng năng suất may, công ty đó may được bao nhiêu chiếc áo?
\choice
{\True $20\,000$ chiếc}
{$18\,000$ chiếc}
{$22\,000$ chiếc}
{$25\,000$ chiếc}
\loigiai{
1 ngày may được: $12\,500 : 5 = 2500$ chiếc.\\
8 ngày may được: $2500 \times 8 = 20\,000$ chiếc.\\
\textbf{Đáp án của câu hỏi là A}
}
\end{ex}

\begin{ex} %[4D1V02]
Tìm số dư trong phép chia $56\,789 : 6$.
\choice
{$3$}
{$4$}
{$1$}
{\True $5$}
\loigiai{
$56\,789 = 6 \times 9464 + 5$. Số dư là 5.\\
\textbf{Đáp án của câu hỏi là D}
}
\end{ex}

\begin{ex} %[4D1V02]
Một người mua $4$ hộp bánh cùng loại hết $96\,000$ đồng. Nếu mua $7$ hộp bánh như thế thì phải trả số tiền là
\choice
{$158\,000$ đồng}
{$172\,000$ đồng}
{\True $168\,000$ đồng}
{$180\,000$ đồng}
\loigiai{
1 hộp bánh giá: $96\,000 : 4 = 24\,000$ đồng.\\
7 hộp bánh giá: $24\,000 \times 7 = 168\,000$ đồng.\\
\textbf{Đáp án của câu hỏi là C}
}
\end{ex}

\begin{ex} %[4D1C02]
Tìm một số biết rằng lấy số đó chia cho 7 rồi cộng với $12\,500$ thì được số tròn chục nghìn bé nhất có 5 chữ số.
\choice
{$0$}
{\True Không có số tự nhiên (hay nếu số tròn chục nghìn là $20\,000 \Rightarrow 52\,500$)}
{$52\,500$}
{$70\,000$}
\loigiai{
Số tròn chục nghìn bé nhất có 5 chữ số là $10\,000$. Khi đó số đó là âm. Nếu số tròn chục nghìn tiếp theo là $20\,000$ thì $(20\,000 - 12\,500) \times 7 = 7500 \times 7 = 52\,500$.\\
\textbf{Đáp án của câu hỏi là B}
}
\end{ex}

\textbf{Phần 2. Trắc nghiệm đúng sai}
\setcounterpageref{ex}{0}

\begin{ex} %[4D1H02]
Xét tính đúng sai của các phép tính sau:
\choiceTF
{\True $35\,000 + 45\,000 = 80\,000$.}
{$60\,000 - 25\,000 = 45\,000$.}
{\True $15\,000 \times 4 = 60\,000$.}
{\True $72\,000 : 8 = 9000$.}
\loigiai{
\begin{itemize}
    \item $35\,000 + 45\,000 = 80\,000$. Đúng.
    \item $60\,000 - 25\,000 = 35\,000 \neq 45\,000$. Sai.
    \item $15\,000 \times 4 = 60\,000$. Đúng.
    \item $72\,000 : 8 = 9000$. Đúng.
\end{itemize}
\begin{itemize}
    \item a) \textbf{Đúng}.
    \item b) \textbf{Sai}.
    \item c) \textbf{Đúng}.
    \item d) \textbf{Đúng}.
\end{itemize}
}
\end{ex}

\begin{ex} %[4D1V02]
Có ba thùng dầu: Thùng 1 chứa $12\,500\text{ l}$, thùng 2 chứa nhiều hơn thùng 1 là $3500\text{ l}$, thùng 3 chứa bằng một nửa thùng 2. Xét tính đúng sai:
\choiceTF
{\True Thùng 2 chứa $16\,000\text{ l}$ dầu.}
{\True Thùng 3 chứa $8000\text{ l}$ dầu.}
{Cả 3 thùng chứa tất cả $35\,000\text{ l}$ dầu.}
{\True Cả 3 thùng chứa tất cả $36\,500\text{ l}$ dầu.}
\loigiai{
\begin{itemize}
    \item Thùng 2: $12\,500 + 3500 = 16\,000\text{ l}$. Đúng.
    \item Thùng 3: $16\,000 : 2 = 8000\text{ l}$. Đúng.
    \item Tổng 3 thùng: $12\,500 + 16\,000 + 8000 = 36\,500\text{ l} \neq 35\,000\text{ l}$. Ý 3 Sai, Ý 4 Đúng.
\end{itemize}
\begin{itemize}
    \item a) \textbf{Đúng}.
    \item b) \textbf{Đúng}.
    \item c) \textbf{Sai}.
    \item d) \textbf{Đúng}.
\end{itemize}
}
\end{ex}

\textbf{Phần 3. Trả lời ngắn}
\setcounterpageref{ex}{0}

\begin{ex} %[4D1H02]
Tính giá trị của biểu thức: $45\,000 + 12\,000 \times 3$.
\shortans{81000}
\loigiai{
\textbf{Đáp số:} $81000$\\
$45\,000 + 36\,000 = 81\,000$.
}
\end{ex}

\begin{ex} %[4D1H02]
Tính hiệu: $90\,000 - 45\,678$.
\shortans{44322}
\loigiai{
\textbf{Đáp số:} $44322$\\
$90\,000 - 45\,678 = 44\,322$.
}
\end{ex}

\begin{ex} %[4D1V02]
Tìm $x$ biết: $x : 6 = 14\,250$.
\shortans{85500}
\loigiai{
\textbf{Đáp số:} $85500$\\
$x = 14\,250 \times 6 = 85\,500$.
}
\end{ex}

\begin{ex} %[4D1V02]
Một nông trường nuôi $15\,600$ con gà, số con vịt ít hơn số con gà $4200$ con. Tính tổng số gà và vịt của nông trường.
\shortans{27000}
\loigiai{
\textbf{Đáp số:} $27000$\\
Số vịt: $15\,600 - 4200 = 11\,400$ con.\\
Tổng: $15\,600 + 11\,400 = 27\,000$ con.
}
\end{ex}

\begin{ex} %[4D1V02]
Một tổ công nhân có 6 người đắp được $480\text{ m}$ đường. Hỏi nếu có 9 người với cùng năng suất thì đắp được bao nhiêu mét đường?
\shortans{720}
\loigiai{
\textbf{Đáp số:} $720$\\
1 người: $480 : 6 = 80\text{ m} \Rightarrow 9$ người: $80 \times 9 = 720\text{ m}$.
}
\end{ex}

\begin{ex} %[4D1C02]
Tính tổng: $12\,345 + 23\,456 + 34\,567 + 45\,678 + 56\,789$.
\shortans{172835}
\loigiai{
\textbf{Đáp số:} $172835$\\
Tổng $= 172\,835$.
}
\end{ex}
